% !TeX root = main.tex
\chapter{Conception et réalisation}

\section{Introduction}

Ce chapitre décrit le passage entre les besoins formalisés dans les parties précédentes et leur traduction concrète dans la plateforme \textit{HireCue}. L'enjeu n'était pas d'ajouter des modules isolés, mais de faire évoluer un système déjà en service en conservant une architecture lisible, des parcours exploitables et une logique métier défendable.

Dans le cadre du stage chez Teligencia Labs, l'accent a surtout porté sur des briques qui devaient apporter une valeur immédiate au recrutement : score explicable, recommandations côté recruteur, roadmap de compétences, chatbot contextuel, questions personnalisées et workflows d'automatisation. La lecture qui suit revient d'abord sur l'architecture générale et sur la modélisation UML, puis sur la conception fonctionnelle, la conception technique et l'implémentation des fonctionnalités effectivement rattachées au projet.

Pour clarifier l'organisation du chapitre, les premières sections décrivent les choix de conception : architecture, modèles UML, rôle fonctionnel des modules et principes techniques retenus. Les sections suivantes présentent ensuite la réalisation concrète, l'intégration décisionnelle et les workflows automatisés, avant de revenir sur les difficultés rencontrées.

\section{Architecture globale du système}

\textit{HireCue} repose sur une architecture distribuée dans laquelle le frontend React \cite{react_docs}, le backend Node.js / Express \cite{node_docs}, la base MySQL \cite{mysql_docs}, les services IA et les workflows externes coopèrent sans se confondre. Le frontend porte les parcours visibles du candidat, du recruteur et de l'administrateur. Il transmet les actions utilisateur au backend via des APIs REST, puis restitue les résultats sous forme de dashboards, de rapports, de roadmaps ou de listes de candidatures exploitables.

Le backend centralise ensuite la logique métier. Il orchestre les appels aux services applicatifs, interroge la base, gère les quotas et coordonne plusieurs traitements qui ne doivent pas être bloqués par une unique requête synchrone. Lorsqu'un flux externe est nécessaire, le backend peut transmettre la demande à \textit{n8n} \cite{n8n_docs}. Ce dernier joue alors le rôle d'orchestrateur pour les imports, les enchaînements de nettoyage, les interactions avec PhantomBuster \cite{phantombuster_docs} et certains flux de communication. Les données consolidées reviennent ensuite vers MySQL, puis peuvent être exploitées par les modules métier ou par les services IA.

Les services IA sont intégrés de manière ciblée. Ils interviennent pour produire une explication, une roadmap, une reformulation ou un contenu structuré, mais ils ne remplacent pas les calculs métier déterministes. Dans la pratique, le flux peut donc se lire ainsi : le frontend React sollicite le backend API, le backend active si besoin un workflow \textit{n8n}, ce workflow dialogue avec PhantomBuster ou avec d'autres services externes, les données sont validées et stockées dans MySQL, puis les modules IA viennent enrichir certains résultats déjà construits.

\begin{figure}[H]
    \centering
    \includegraphics[width=0.95\textwidth]{\detokenize{Diagramme d’architecture.png}}
    \caption{Architecture globale du système HireCue}
    \label{fig:architecture-globale-hirecue}
\end{figure}

\section{Conception UML}

La modélisation UML a servi de support pour clarifier l'organisation du système avant de décrire les détails d'implémentation. Dans un projet comme \textit{HireCue}, cette étape est utile parce qu'elle permet de visualiser les liens entre les entités, les interactions entre composants et le déroulement de plusieurs traitements métier sans basculer trop tôt dans le niveau du code.

\subsection{Diagramme de classes}

Le diagramme de classes montre que le noyau de \textit{HireCue} reste structuré autour de quelques entités pivot. La table User porte les informations de profil, de rôle, de rattachement à l'entreprise et de quotas IA. La table Job décrit les offres, tandis que CandidateJob joue un rôle central en reliant candidat, poste, résultats de tests et état de la candidature. C'est d'ailleurs autour de cette entité que se greffent le score explicable, la roadmap, certaines alertes de risque, les recommandations et le réengagement.

Cette organisation confirme que les modules réalisés pendant le stage ne sont pas des briques ajoutées à côté du système. Ils réutilisent les données métier existantes, puis les prolongent par de la persistance, du cache ou des journaux de traitement.

\begin{figure}[H]
    \centering
    \includegraphics[width=0.95\textwidth]{\detokenize{diagramme de classe.png}}
    \caption{Diagramme de classes de la plateforme HireCue}
    \label{fig:diagramme-classes-hirecue}
\end{figure}

\subsection{Diagrammes de cas d'utilisation}

Les diagrammes de cas d'utilisation aident à replacer les fonctionnalités dans le point de vue des acteurs. Dans \textit{HireCue}, ils permettent surtout de montrer que les modules intelligents ne sont pas présentés comme des traitements abstraits, mais comme des services rattachés à des actions concrètes du candidat ou du recruteur.

\paragraph{Cas d'utilisation global.}
Le diagramme global met principalement en scène deux acteurs : le candidat et le recruteur. Le candidat y apparaît comme l'utilisateur qui consulte ses notifications, téléverse son CV, passe des tests, suit ses candidatures, consulte les détails de son score, visualise sa feuille de route et parcourt la liste des emplois. Deux extensions attirent l'attention : l'accès au rapport détaillé et l'usage du chatbot de la roadmap. Le recruteur intervient de son côté sur la consultation des alertes de risque et sur la modération des emplois. Ce schéma reste volontairement synthétique. Il n'explicite pas toutes les briques internes d'IA, mais il montre bien les parcours qui justifient ensuite le score explicable, la roadmap et les alertes.

\begin{figure}[H]
    \centering
    \includegraphics[width=0.95\textwidth]{Casdutilisation_globale.png}
    \caption{Diagramme de cas d'utilisation global de HireCue}
    \label{fig:cas-utilisation-global-hirecue}
\end{figure}

\paragraph{Cas d'utilisation raffiné.}
Le diagramme raffiné se concentre sur une vue recruteur plus détaillée autour de la modération des emplois. Il montre que cette action ne se limite pas à afficher une liste. Le recruteur doit d'abord s'authentifier, consulter la liste des candidats, la filtrer, puis ouvrir le détail d'un profil. À partir de ce point, plusieurs extensions deviennent possibles : voir le rapport, voir les recommandations IA ou gérer le réengagement. Le réengagement lui-même inclut la consultation du profil et l'envoi d'un message. Cette vue est utile, car elle relie dans un même scénario la lecture du score, l'aide à la décision et l'action de relance.

\begin{figure}[H]
    \centering
    \includegraphics[width=0.95\textwidth]{casdutilisation_raffine.png}
    \caption{Diagramme de cas d'utilisation raffiné de HireCue}
    \label{fig:cas-utilisation-raffine-hirecue}
\end{figure}

\subsection{Diagrammes de cas d'utilisation n8n}

Les diagrammes suivants complètent la modélisation UML en décrivant les principaux workflows \textit{n8n} reliés à LinkedIn. Ils montrent comment les actions initiées par un administrateur ou un recruteur sont relayées par le backend, orchestrées par \textit{n8n}, puis exécutées à travers UniPile ou les services LinkedIn. Ces représentations permettent aussi de distinguer les flux de publication, les flux de suivi d'engagement et les flux de prospection automatisée.

\subsubsection{Publication et suivi automatique de l'engagement des posts HireCue}

Le premier diagramme représente un workflow de publication de posts \textit{HireCue} sur LinkedIn, complété par un mécanisme de suivi automatique des interactions. L'administrateur ou le recruteur intervient comme acteur métier à l'origine du besoin de publication, tandis que le backend \textit{HireCue} prépare les données nécessaires et transmet l'exécution au workflow \textit{n8n}. Le schéma montre ensuite la préparation du payload LinkedIn, la publication du post via UniPile, l'enregistrement de l'identifiant du post LinkedIn et la mise à jour du statut de publication, qui peut être marqué comme publié ou en échec.

La seconde partie du diagramme décrit l'\textit{interaction tracking}. Après la publication, \textit{n8n} récupère les posts publiés, obtient leurs détails et leur identifiant social, puis suit les commentaires et les réactions des utilisateurs LinkedIn. Ces interactions sont sauvegardées dans la base MySQL afin de conserver une trace exploitable. Le workflow calcule ensuite des statistiques d'engagement qui peuvent être affichées dans un tableau de bord administrateur. Cette lecture montre que la publication n'est pas une action isolée : elle s'inscrit dans une boucle de suivi qui permet d'analyser les commentaires, les réactions, les mentions J'aime et les performances globales du contenu.

\begin{figure}[H]
    \centering
    \includegraphics[width=0.95\textwidth]{usecasenormalpost.png}
    \caption{Diagramme de cas d'utilisation n8n de publication et suivi des posts HireCue}
    \label{fig:cas-utilisation-n8n-posts-hirecue}
\end{figure}

\subsubsection{Publication des offres d'emploi sur LinkedIn}

Le deuxième diagramme se concentre sur le workflow de publication d'une offre d'emploi \textit{HireCue} vers LinkedIn. Le recruteur ou l'administrateur commence par sélectionner une offre interne. Le contenu de l'offre est ensuite transformé en post LinkedIn, puis enrichi par un lien de candidature \textit{HireCue}. Cette étape garantit que la visibilité LinkedIn reste reliée au parcours de candidature géré par la plateforme.

Le diagramme met aussi en évidence le rôle de \textit{n8n} et d'UniPile dans l'exécution de la publication. \textit{n8n} orchestre l'appel de publication, tandis qu'UniPile joue le rôle d'interface API avec LinkedIn. Une fois l'offre publiée, le workflow peut suivre les réactions et les commentaires des candidats, afficher les statistiques du post et, lorsque cela est nécessaire, s'appuyer sur un service IA ou LLM pour générer une réponse contenant le lien de candidature. Dans la logique applicative de \textit{HireCue}, le backend conserve le statut de publication, l'identifiant du post et les informations utiles au suivi, afin que le recruteur puisse retrouver l'état du flux sans dépendre uniquement de LinkedIn.

\begin{figure}[H]
    \centering
    \includegraphics[width=0.95\textwidth]{usecaseoffrejob.png}
    \caption{Diagramme de cas d'utilisation n8n de publication des offres d'emploi sur LinkedIn}
    \label{fig:cas-utilisation-n8n-offres-linkedin}
\end{figure}

\subsubsection{Automatisation des messages LinkedIn pour la prise de contact avec des PDG et fondateurs}

Le troisième diagramme présente un workflow d'outreach LinkedIn destiné aux CEO, founders et owners. Il est organisé en trois parties complémentaires. La première, consacrée à la détection des profils cibles, commence par la planification d'une recherche LinkedIn par le recruteur ou l'administrateur. Les profils de CEO ou de fondateurs sont ensuite recherchés via UniPile, normalisés, filtrés selon leur rôle de décisionnaire, puis stockés dans le backend.

La deuxième partie porte sur l'envoi des demandes de connexion. Les leads à contacter sont récupérés, une note de connexion personnalisée est générée, puis la demande LinkedIn est envoyée via UniPile. Le workflow marque ensuite l'invitation comme envoyée, ce qui permet de maintenir une trace claire de chaque lead. La troisième partie intervient après acceptation : le workflow détecte la connexion acceptée, prépare un message de suivi, l'envoie automatiquement et met à jour le statut d'outreach. Le diagramme associe enfin ces statuts à une zone d'analyse et de suivi, comprenant le calcul des KPI, le nombre d'invitations envoyées, le taux d'acceptation et le suivi des messages de relance. Cette organisation montre que \textit{n8n} et UniPile ne servent pas uniquement à envoyer des messages, mais aussi à maintenir une traçabilité exploitable dans le backend et dans la base MySQL.

\begin{figure}[H]
    \centering
    \includegraphics[width=0.95\textwidth]{usecaseceooutrich.png}
    \caption{Diagramme de cas d'utilisation n8n d'automatisation des messages LinkedIn}
    \label{fig:cas-utilisation-n8n-outreach-ceo}
\end{figure}

\subsection{Diagramme de séquence}

Les diagrammes de séquence rendent plus lisible la coopération entre interface, backend, base de données, cache et services IA. Dans \textit{HireCue}, ils montrent bien que la génération d'un résultat ne repose pas sur un seul appel à un modèle, mais sur une chaîne d'étapes qui comprend collecte des données, vérification de cache, calcul métier, validation de format et persistance.

\subsubsection{Diagramme de séquence - Génération de roadmap}

Le premier scénario décrit la génération d'une roadmap depuis l'ouverture de la page candidate. Le frontend transmet une requête au backend, qui charge le profil du candidat, le poste cible et les compétences utiles. Le \textit{RoadmapService} calcule ensuite les \textit{skill gaps}, puis vérifie si une roadmap équivalente est déjà disponible en cache. Lorsque le résultat existe déjà, il peut être retourné sans régénération, ce qui limite les coûts et améliore le temps de réponse.

En l'absence de cache, le service envoie un contexte structuré au service IA avec les compétences manquantes. La réponse n'est pas acceptée aveuglément. Le diagramme montre une étape de validation JSON, puis un mécanisme de reprise si le format reste invalide. Si la seconde tentative échoue encore, le système bascule vers un fallback déterministe. Cette logique est importante dans le projet, car elle montre que la génération n'est pas laissée sans contrôle et qu'une sortie exploitable reste attendue même en cas d'échec du modèle.

\begin{figure}[H]
    \centering
    \includegraphics[width=0.95\textwidth]{\detokenize{Diagramme de séquence — Génération de roadmap.png}}
    \caption{Diagramme de séquence de génération de roadmap}
    \label{fig:sequence-roadmap}
\end{figure}

\subsubsection{Diagramme de séquence - Recommandations IA}

Le second diagramme commence côté recruteur, au moment de l'ouverture de la page des recommandations. Le backend interroge d'abord le \textit{RecommendationService}, qui vérifie l'existence d'un résultat en cache. Si un jeu de recommandations compatible est déjà présent, le système retourne directement la réponse au frontend. Dans le cas contraire, le service charge le job et les candidats compatibles depuis la base, compare les profils et calcule un classement ainsi qu'un \textit{fit band}.

Le point essentiel de ce diagramme est l'ordre des étapes. La comparaison candidat/job, le classement et la détection des signaux de risque apparaissent avant l'appel au service IA. Le LLM intervient ensuite pour enrichir l'explication retournée au recruteur, mais il ne décide pas lui-même de l'ordre final. Cette lecture est cohérente avec le positionnement \textit{rules-first} retenu dans le projet : le résultat principal reste déterministe, tandis que la couche IA sert à rendre l'explication plus lisible.

\begin{figure}[H]
    \centering
    \includegraphics[width=0.95\textwidth]{\detokenize{Diagramme de séquence — Recommandations IA.png}}
    \caption{Diagramme de séquence des recommandations IA}
    \label{fig:sequence-recommandations-ia}
\end{figure}

\subsubsection{Diagramme de séquence - Score explicable}

Le troisième scénario concerne la consultation d'un rapport de score par un candidat ou un recruteur. Le frontend appelle l'API de score, puis le backend récupère dans la base les données reliées à la candidature : candidate\_job, job, CV, tests et scores. Le \textit{ScoreService} calcule alors le score explicable en vérifiant d'abord la disponibilité du final\_tests\_score. Lorsque ce score existe, il est utilisé directement. Sinon, le service regarde si des poids de calcul sont configurés et applique soit une moyenne pondérée, soit une moyenne plus simple.

Une fois le score consolidé, une explication structurée est générée, puis sauvegardée dans score\_reports. Ce diagramme confirme que l'explication intervient après le calcul du score et non à sa place. Il rappelle aussi qu'un même endpoint doit produire un résultat suffisamment stable pour être relu ensuite, exporté ou consulté depuis des parcours différents.

\begin{figure}[H]
    \centering
    \includegraphics[width=0.95\textwidth]{\detokenize{Diagramme de séquence — Score explicable.png}}
    \caption{Diagramme de séquence du score explicable}
    \label{fig:sequence-score-explicable}
\end{figure}


\section{Conception fonctionnelle}

La conception fonctionnelle de \textit{HireCue} décrit le rôle attendu de chaque module dans les parcours utilisateur, sans détailler encore son implémentation. Chaque brique a été rattachée à un besoin concret, soit pour mieux expliquer un résultat, soit pour assister le recruteur sans lui retirer la décision, soit pour rendre certaines opérations répétitives plus fluides.

\subsection{Gestion des candidats}

La gestion des candidats constitue le socle sur lequel reposent la plupart des autres modules. Elle couvre le profil utilisateur, le CV, les informations personnelles, les compétences, les résultats de tests, les candidatures rattachées à un poste et les états de suivi. Dans le projet, cette couche ne sert pas uniquement à conserver des informations administratives. Elle alimente directement les traitements de score, la comparaison avec un poste cible, la roadmap et le chatbot.

Cette centralisation est importante pour les deux principaux profils utilisateur. Le candidat doit retrouver un suivi lisible de sa candidature et des retours associés. Le recruteur, de son côté, a besoin d'un accès cohérent aux informations du profil pour analyser un dossier, lancer des recommandations ou préparer un entretien sans naviguer entre plusieurs sources dispersées.

\subsection{Rapport de score explicable}

Le rapport de score explicable répond à un besoin essentiel : permettre aux candidats et aux recruteurs de comprendre la signification d'un score. Un résultat numérique seul ne suffit pas ; il doit être accompagné d'une explication claire sur les éléments qui ont influencé son calcul.

Cette fonctionnalité vise à rendre le processus d'évaluation plus transparent en présentant les différentes composantes du score, telles que les compétences, les résultats des tests ou d'autres critères pris en compte par la plateforme. Elle contribue ainsi à renforcer la confiance des utilisateurs dans les résultats obtenus.

La restitution ne se limite pas à une note globale. Elle met en évidence les sous-scores, les points forts, les axes d'amélioration et les recommandations éventuelles. Pour le candidat, elle constitue un retour constructif sur son profil, tandis que pour le recruteur, elle fournit une aide à la décision plus claire et plus facilement interprétable.

\subsection{Recommandations IA côté recruteur}

Ce module a pour objectif d'aider les recruteurs à identifier rapidement les candidats les plus pertinents pour une offre, tout en conservant la maîtrise de la décision finale. Il ne remplace pas l'évaluation humaine, mais fournit des indicateurs utiles tels qu'un score d'adéquation, les compétences correspondantes, les écarts identifiés et une explication associée.

Dans \textit{HireCue}, les recommandations sont intégrées aux autres fonctionnalités de la plateforme. Elles peuvent être complétées par l'accès au score explicable, aux questions personnalisées ou aux mécanismes de réengagement des candidats.

Afin de garantir une présentation claire, le système met en avant un nombre limité de profils recommandés, généralement les trois candidats les plus adaptés à une offre. Chaque recommandation est accompagnée d'une justification permettant au recruteur de comprendre les raisons de sa sélection et de poursuivre le processus de recrutement à travers des actions adaptées.
\subsection{Skill Gap Roadmap}

La roadmap de compétences a été conçue pour aider les candidats à comprendre comment progresser après une évaluation. Au-delà d'un simple résultat, elle transforme les écarts identifiés entre un profil et un poste en un plan d'amélioration structuré.

Cette fonctionnalité apporte une dimension d'accompagnement à \textit{HireCue}. Au lieu de se limiter à une évaluation, la plateforme fournit des recommandations permettant au candidat d'identifier les compétences à développer et les priorités de progression.

La roadmap est générée à partir d'une candidature et d'un poste cible. Elle met en évidence les compétences manquantes, les axes d'amélioration et les actions recommandées pour renforcer l'adéquation du profil au poste. Un export PDF peut également être généré afin de faciliter la consultation et le partage du plan de progression.

\subsection{Chatbot RAG}

Le chatbot lié à la roadmap intervient comme un prolongement de cette dernière. Son objectif est de répondre aux questions du candidat en s'appuyant sur le contexte déjà disponible : contenu de la roadmap, compétences observées, poste cible et historique des échanges. L'intérêt fonctionnel de cette brique est de garder la discussion ancrée dans un périmètre utile au lieu d'ouvrir une conversation trop générale.

Dans le projet, ce module apporte surtout de la fluidité. Il permet au candidat d'approfondir un conseil, de demander une reformulation ou de clarifier une étape sans devoir sortir du parcours. Son utilité dépend donc autant de la qualité du contexte transmis que de la génération elle-même.

\subsection{Questions personnalisées}

Les questions personnalisées s'adressent au recruteur et prolongent l'analyse du profil vers la préparation de l'entretien. Le besoin fonctionnel est clair : une liste de questions standard ne permet pas toujours d'explorer les points faibles ou les zones d'incertitude d'un candidat. En s'appuyant sur le poste, sur les écarts de compétences et sur les recommandations déjà calculées, le module propose des questions plus ciblées.

Cette brique renforce le caractère opérationnel de \textit{HireCue}. Elle traduit un résultat analytique en un usage concret pour la suite du recrutement, tout en gardant un format assez simple pour être directement affiché dans le dashboard recruteur.

\subsection{Candidate Engagement \& Motivation Insights}

Le suivi d'engagement candidat vise à donner plus de lisibilité au parcours après l'évaluation. Dans la structure du système, cette brique s'appuie sur des marqueurs de progression et sur des informations utiles à la lecture du comportement du candidat face à sa roadmap, à ses notifications ou à certaines relances. L'objectif n'est pas de fabriquer artificiellement une motivation, mais de repérer si le parcours reste actif, incomplet ou en perte d'interaction.

Du point de vue métier, cette fonctionnalité peut aider à mieux interpréter certains cas limites. Un candidat qui consulte sa roadmap, revient sur son parcours et interagit avec les modules d'accompagnement n'est pas dans la même situation qu'un profil qui reste inactif. Cette distinction nourrit ensuite d'autres briques, notamment les alertes ou le réengagement.

\subsection{Réengagement des candidats}

Le réengagement vise à mieux exploiter le vivier déjà présent dans la plateforme. Plutôt que de relancer exclusivement une recherche externe, le recruteur peut retrouver des profils déjà connus, identifier ceux qui restent pertinents pour un poste et lancer une relance structurée. Ce module concerne d'abord le dashboard recruteur, mais il a aussi un effet sur l'expérience candidat lorsqu'une nouvelle opportunité lui est proposée.

Dans le projet, cette brique apporte une valeur métier directe. Elle aide à transformer des candidatures anciennes ou peu actives en un vivier mieux valorisé, sans imposer une reprise manuelle de toute l'analyse.

\subsection{Relance admin des candidats inactifs}

Au-delà du réengagement piloté par le recruteur, un besoin plus transversal concerne la relance des candidats inactifs depuis le dashboard Admin. L'objectif est de disposer d'un déclencheur simple permettant d'envoyer un email à des profils restés sans activité depuis une durée voisine de soixante jours. Cette fonctionnalité vise moins la sélection pour un poste précis que la réactivation générale du vivier présent sur la plateforme.

Dans le périmètre actuel, cette brique doit être décrite avec prudence. Selon le niveau de finalisation effectivement retenu, elle peut correspondre à un bouton déjà intégré, à un mécanisme documenté ou à une fonctionnalité en cours de consolidation. Dans tous les cas, sa place dans le rapport est justifiée, car elle prolonge la logique d'engagement candidat et complète les modules de relance déjà présents côté recruteur.

\subsection{Smart Candidate Risk Alerts}

Les alertes de risque ont pour rôle d'attirer l'attention du recruteur sur des dossiers qui demandent une lecture plus fine. Il peut s'agir d'un engagement faible, d'informations manquantes, d'un score bas ou d'une incohérence entre plusieurs signaux. L'objectif fonctionnel n'est pas de disqualifier automatiquement un candidat, mais de faire remonter des points de vigilance qui méritent une vérification humaine.

Dans \textit{HireCue}, cette fonctionnalité renforce l'aide à la décision sans la remplacer. Elle s'insère naturellement dans le dashboard recruteur, aux côtés des recommandations et du score détaillé.

\subsection{Scraping des données externes}

Le scraping des données externes permet à \textit{HireCue} d'enrichir sa base d'informations à partir de sources telles que LinkedIn. Son objectif est de faciliter le sourcing de candidats et la collecte d'offres d'emploi afin de compléter les données déjà présentes dans la plateforme.

Cette fonctionnalité est principalement utilisée pour rechercher de nouveaux profils candidats et récupérer des offres externes. Les données collectées ne sont pas exploitées directement ; elles sont d'abord vérifiées, nettoyées et normalisées afin de garantir leur cohérence avec le modèle de données interne.

Sur le plan technique, PhantomBuster est utilisé pour l'extraction des données, tandis que \textit{n8n} orchestre les différents workflows. Les informations récupérées sont ensuite traitées, dédupliquées et transformées avant leur intégration dans la base de données ou leur affichage dans les tableaux de bord de la plateforme.

Ainsi, la valeur de cette fonctionnalité ne réside pas uniquement dans la collecte de données externes, mais surtout dans leur transformation en informations fiables et exploitables pour les processus de recrutement de \textit{HireCue}.
\section{Conception technique}

La conception technique retenue pendant le stage cherche à séparer clairement ce qui relève de l'interface, de la logique métier, de la persistance, des traitements intelligents et des intégrations externes. Cette séparation était nécessaire, car les nouvelles fonctionnalités ne pouvaient pas être ajoutées proprement si tout le système était concentré dans une seule couche.

\subsection{Architecture backend}

L'architecture backend de \textit{HireCue} repose sur une organisation en couches composée des routes API, des contrôleurs, des services métier, des repositories et des modèles Sequelize. Cette structure favorise la modularité, la maintenabilité et la séparation des responsabilités.

Les routes exposent les services consommés par le frontend, tandis que les contrôleurs gèrent les requêtes, les validations et l'authentification avant de déléguer les traitements aux services métier. Ces derniers regroupent l'essentiel de la logique applicative, notamment le calcul des scores, la génération des rapports explicables, des roadmaps de compétences, des recommandations intelligentes et des fonctionnalités de réengagement.

Pour les modules d'intelligence artificielle, le backend reconstruit le contexte métier à partir des données stockées dans MySQL avant d'interroger les services IA. Les résultats générés sont ensuite validés, éventuellement mis en cache, puis renvoyés au frontend dans un format exploitable.

L'accès aux données est assuré par les repositories et les modèles Sequelize, qui encapsulent les opérations réalisées sur la base MySQL. Cette approche améliore la cohérence des accès aux données et facilite l'évolution du système.

Grâce à cette architecture, les différentes fonctionnalités de \textit{HireCue} peuvent être intégrées de manière cohérente tout en conservant de bonnes performances, une maintenance simplifiée et une meilleure évolutivité. Elle facilite également la gestion des traitements asynchrones, de la mise en cache et du suivi des usages des services d'intelligence artificielle.

\begin{figure}[H]
    \centering
    \includegraphics[width=\textwidth]{archi_backend.png}
    \caption{Architecture backend de la plateforme HireCue}
\end{figure}

\subsection{Base de données}

La base de données MySQL occupe une position centrale dans l'architecture de \textit{HireCue}. Elle assure non seulement la persistance des données métier classiques, mais également le stockage des résultats de calcul, des historiques d'interaction, des caches applicatifs et de certaines traces liées aux fonctionnalités intelligentes.
Parmi les tables les plus structurantes, la table users regroupe les informations relatives aux utilisateurs de la plateforme, notamment leurs profils et leurs rôles. La table joblist centralise les informations associées aux offres d'emploi publiées. La table candidate\_job matérialise la relation entre un candidat et une offre donnée. La table cv\_info complète cet ensemble en stockant les informations issues du CV ou associées au profil candidat.

Dans le cadre des fonctionnalités développées dans ce projet, plusieurs tables spécialisées interviennent également de manière significative. La table score\_reports permet de persister les rapports explicables générés à partir des scores candidats ou recruteurs. La table candidate\_skill\_roadmaps est utilisée pour stocker les roadmaps de progression produites à partir de l'analyse des écarts de compétences.

La base de données joue également un rôle dans le pilotage des fonctionnalités génératives ou assistées par intelligence artificielle. À ce titre, la table LLMConsumptions permet de suivre la consommation associée à certains appels IA, contribuant ainsi à la traçabilité et à la régulation de leur usage. D'autres tables, liées aux notifications, aux rappels ou aux journaux d'activité, participent également au bon fonctionnement de fonctionnalités telles que le réengagement de candidats ou l'envoi d'emails de relance après une période d'inactivité.
Cette centralité de la base de données contribue directement à la cohérence du backend et à la stabilité des fonctionnalités avancées développées dans le cadre de la plateforme.

\begin{figure}[H]
    \centering
    \includegraphics[width=\textwidth]{archi_basededonnes.png}
    \caption{Architecture de la base de données de HireCue}
\end{figure}


\subsection{Intégration n8n}

\textit{n8n} intervient comme une couche d'orchestration séparée du backend principal. Dans le cadre du projet, ce choix a été retenu afin d'éviter que toute la logique de scraping, de transformation et de communication avec des services tiers soit mélangée aux routes métier classiques. Cette séparation simplifie les diagnostics, rend les workflows plus faciles à faire évoluer et limite l'impact d'un échec externe sur le cœur applicatif.

Concrètement, \textit{n8n} est mobilisé pour lancer des recherches, recevoir des données extraites, normaliser certains champs, dédoublonner les résultats et transmettre au backend des structures exploitables. Il peut aussi servir de point d'orchestration pour des scénarios de communication ou de suivi d'exécution.

Les échanges entre le backend et \textit{n8n} reposent principalement sur des webhooks et des appels HTTP. Lorsqu'un utilisateur déclenche une action depuis l'interface, le backend prépare les paramètres utiles, par exemple un mot-clé de recherche, une localisation, un identifiant de job ou un statut de campagne, puis appelle le webhook du workflow. \textit{n8n} exécute ensuite les étapes prévues et renvoie au backend un résultat normalisé ou un état d'exécution. Cette boucle permet de garder le backend comme point de contrôle métier, tout en laissant à \textit{n8n} l'orchestration des connecteurs externes.

\begin{figure}[H]
    \centering
    \includegraphics[width=0.35\textwidth]{N8n-logo-new.svg.png}
    \caption{Logo de n8n, outil d'orchestration des workflows}
    \label{fig:n8n-logo}
\end{figure}

\subsection{Intégration PhantomBuster}

PhantomBuster \cite{phantombuster_docs} est utilisé lorsque la plateforme doit interagir avec des données externes LinkedIn. Son rôle est de lancer des agents capables d'extraire des profils, des emails professionnels lorsque ceux-ci sont disponibles, ou encore des offres d'emploi. L'outil n'est toutefois pas intégré seul. Les résultats passent ensuite par \textit{n8n}, qui assure les étapes de normalisation, de déduplication et de préparation avant stockage ou import via l'API backend.

Cette intégration exige une certaine prudence, car les données récupérées ne sont ni homogènes ni garanties. Une partie du travail a donc consisté à traiter PhantomBuster comme une source externe utile, mais non suffisante en elle-même. La qualité finale dépend de la chaîne complète : extraction, nettoyage, vérification, mapping et insertion.

\begin{figure}[H]
    \centering
    \includegraphics[width=0.35\textwidth]{PhantomBuster.png}
    \caption{Logo de PhantomBuster utilisé pour l'extraction de données externes}
    \label{fig:phantombuster-logo}
\end{figure}

\subsection{Intégration UniPile}

UniPile \cite{unipile_docs} est utilisé dans \textit{HireCue} pour gérer les interactions avec LinkedIn, notamment l'envoi de messages, le suivi des conversations et la publication de contenus ou d'offres d'emploi. Les actions sont préparées par le backend puis exécutées via l'API UniPile, éventuellement à travers des workflows \textit{n8n}.

Cette solution complète PhantomBuster : tandis que PhantomBuster est dédié au sourcing et à l'extraction de données, UniPile prend en charge les échanges et les publications sur LinkedIn. Cette séparation permet d'utiliser chaque outil selon son domaine de spécialisation.


\begin{figure}[H]
    \centering
    \includegraphics[width=0.45\textwidth]{logo-unipile-fond-fonce-1-2-1-1-1500x679.png}
    \caption{Logo de UniPile utilisé pour les intégrations LinkedIn}
    \label{fig:logo-unipile}
\end{figure}

\subsection{Gestion du cache et des traitements asynchrones}

Le cache est présent pour éviter de régénérer inutilement des résultats coûteux, en particulier pour les score\_reports, les roadmaps et certaines recommandations. Cette logique est importante dans \textit{HireCue}, car plusieurs modules combinent des données métier, des quotas IA et des appels qui ne doivent pas être rejoués à chaque consultation si le contexte n'a pas changé.

En parallèle, RabbitMQ permet de déplacer certaines tâches lourdes hors du chemin synchrone. Les recommandations IA et certains jobs de réengagement en sont de bons exemples. Le backend peut enregistrer une demande, la confier à une queue, laisser un worker produire le résultat en arrière-plan, puis laisser le frontend interroger l'état du traitement tant que le calcul n'est pas terminé. Cette approche réduit la charge perçue par l'utilisateur et donne plus de robustesse à des traitements qui pourraient sinon bloquer l'interface ou le serveur.
\section{Implémentation des fonctionnalités}

À partir de cette section, la lecture bascule vers la réalisation. L'objectif n'est pas de reproduire la documentation interne des services, mais d'expliquer comment les choix de conception ont été traduits en modules exploitables dans les parcours candidat et recruteur.

\subsection{Score explicable}

Le score explicable permet de rendre les résultats d'évaluation plus compréhensibles pour les candidats et les recruteurs. Il s'appuie sur plusieurs données de la plateforme, notamment les sous-scores \texttt{Cscore}, \texttt{Pscore}, \texttt{Ptscore} et \texttt{Tscore}, afin de calculer un score global selon des règles métier prédéfinies.

Une fois le score calculé, le système génère une explication détaillant les principaux facteurs qui ont influencé le résultat. Ces informations sont enregistrées dans la table \texttt{score\_reports} afin d'éviter des recalculs inutiles et de garantir la traçabilité des analyses.

La restitution met en évidence les points forts, les axes d'amélioration et les sous-scores associés. Elle offre ainsi un retour constructif au candidat et une aide à la décision plus claire pour le recruteur. Un export PDF du rapport peut également être généré afin de faciliter la consultation et le partage des résultats.


\subsection{AI Recommendations}

Le module de recommandations aide les recruteurs à identifier rapidement les candidats les plus adaptés à une offre. Il repose sur une analyse des compétences correspondantes, des compétences manquantes et du score d'adéquation entre le profil du candidat et le poste.

Le classement des candidats est réalisé par la logique métier de la plateforme afin de garantir des résultats cohérents et explicables. L'intelligence artificielle intervient uniquement pour enrichir les résultats en générant des explications et des synthèses plus faciles à interpréter.

Sur le plan technique, le backend analyse les données du poste et des candidats, calcule les indicateurs de correspondance puis transmet un contexte structuré au service IA. Les résultats peuvent être mis en cache ou traités de manière asynchrone afin d'améliorer les performances.

Cette fonctionnalité ne se limite pas à recommander des profils. Elle permet également d'accéder à d'autres actions, comme la consultation du score explicable, la génération de questions personnalisées ou le lancement de scénarios de réengagement. Elle constitue ainsi un élément central du tableau de bord recruteur.

\subsection{Skill Gap Roadmap}

La \textit{Skill Gap Roadmap} aide le candidat à comprendre les écarts entre son profil et les compétences attendues pour un poste cible. Elle s'appuie sur les données du candidat, son CV, l'offre visée et les compétences requises.

Le système compare d'abord les compétences disponibles avec celles demandées par le poste afin d'identifier les compétences maîtrisées et celles à développer. Cette étape reste déterministe pour garantir un résultat fiable et cohérent.

Après cette analyse, l'intelligence artificielle intervient pour générer une roadmap personnalisée. Celle-ci propose des étapes de progression, des priorités d'apprentissage et des recommandations pratiques adaptées au profil du candidat.

Les résultats sont validés, puis enregistrés dans la table \texttt{candidate\_skill\_roadmaps}. La roadmap est ensuite affichée dans l'interface et peut être exportée au format PDF afin de faciliter son suivi et son partage.

\subsection{Roadmap Chatbot RAG}

Le chatbot associé à la roadmap agit comme un assistant candidat contextualisé. Il s'appuie sur la roadmap générée, le poste cible, les compétences du candidat et l'historique récent des échanges. Cette approche reprend le principe du \textit{Retrieval-Augmented Generation} \cite{lewis_rag}, qui combine récupération de contexte et génération de réponse.

L'objectif du chatbot n'est pas de proposer une conversation générale, mais d'aider le candidat à comprendre sa roadmap, ses compétences manquantes et les étapes de progression recommandées. Les réponses restent donc liées aux données déjà présentes dans \textit{HireCue}.

Lorsqu'une question est posée, le backend récupère la roadmap sauvegardée et une partie de l'historique de conversation, puis transmet ce contexte au service IA. La réponse générée est ensuite enregistrée dans \texttt{candidate\_roadmap\_chat\_messages}, ce qui permet de conserver une continuité entre les échanges tout en limitant la taille du contexte utilisé.

\subsection{Questions personnalisées}

Les questions personnalisées sont générées côté recruteur à partir du triplet candidat, job et analyse existante. Le module peut s'appuyer sur les recommandations, sur les écarts de compétences et sur certains éléments du profil pour proposer des questions plus utiles qu'une trame d'entretien standard.

Le prompt doit rester fortement contraint : format JSON, questions courtes, pas d'hallucination sur des expériences inexistantes et un cadrage assez stable pour que le frontend puisse afficher le résultat sans post-traitement complexe. Ce point est important, car une question pertinente en lecture humaine peut devenir peu exploitable si son format varie trop d'une exécution à l'autre.

\subsection{Candidate Engagement \& Motivation Insights}

Le suivi d'engagement candidat s'inscrit comme un prolongement de la roadmap et du parcours d'accompagnement. L'idée n'est pas de sur-promettre un mécanisme de motivation artificiel, mais de structurer des indicateurs qui aident à comprendre si le candidat consulte son parcours, revient sur ses recommandations ou reste totalement inactif.

Dans sa forme actuelle, cette brique sert surtout à rendre la progression plus lisible et à alimenter d'autres modules, notamment les alertes et le réengagement. Lorsque des marqueurs de progression sont conservés dans l'interface ou dans des structures de cache, ils doivent rester interprétés avec prudence. Le système peut signaler une dynamique de parcours, mais il ne doit jamais inventer une expérience ou une compétence non observée.

\subsection{Re-Engagement Pool}

Le \textit{Re-Engagement Pool} cible les profils déjà connus de la plateforme que le recruteur peut souhaiter relancer pour un poste. Le calcul de ce vivier s'appuie sur des informations déjà présentes dans \textit{HireCue} et peut être délégué à un job asynchrone afin de ne pas bloquer le backend. Cette approche est visible dans les structures de type reengagement\_pool\_jobs et dans le suivi d'état côté interface.

Dans le backlog projet, la logique vise en particulier les candidats restés inactifs pendant une période significative, avec un repère de 60 jours retenu pour orienter la relance. Le module actuellement observé fournit déjà la base technique du réengagement ; l'affinage fin des critères de ciblage s'inscrit davantage dans la consolidation fonctionnelle du dispositif.

\subsection{Scraping des profils LinkedIn}

Le scraping de profils LinkedIn s'appuie sur PhantomBuster et sur \textit{n8n}, notamment dans des scénarios de type \textit{LinkedIn Search to Emails}. Le workflow commence par une recherche externe, puis récupère des profils et, lorsque cela est possible, des emails professionnels. Les données obtenues sont ensuite normalisées, dédoublonnées et préparées avant d'être exposées dans une vue recruteur ou stockées dans un support intermédiaire.

La valeur de cette brique dépend moins du volume brut collecté que de la qualité de l'import final. Le nettoyage des champs, la suppression des doublons et l'alignement avec le modèle interne sont donc des parties centrales de l'implémentation.

\subsection{Workflow CEO / Founder / Owner scraping + email sender}

Dans le périmètre observé, les workflows de prospection de leads de type CEO, founder ou owner apparaissent surtout comme une extension des automatisations \textit{n8n} / PhantomBuster. Leur principe est de récupérer des leads cibles, de vérifier les informations de contact disponibles, puis de préparer un envoi d'emails personnalisés. Lorsqu'un tel scénario est activé, il peut s'appuyer sur Gmail pour l'envoi, sur un suivi de statut et sur un résumé d'exécution remonté vers un canal interne comme Slack.

Cette partie doit toutefois être présentée avec mesure. Elle prolonge l'écosystème d'automatisation autour de \textit{HireCue}, mais elle n'occupe pas la même place structurante que le score, la roadmap ou les recommandations. Dans le rapport, il est donc plus juste de la traiter comme un workflow complémentaire en cours de consolidation.

\subsection{Pipeline d'import des offres LinkedIn}

Le pipeline d'import des offres LinkedIn commence par l'extraction de données externes, suivie d'une normalisation et d'un mapping vers la structure de joblist. Le backend intervient ensuite pour valider les informations reçues, vérifier les doublons à partir d'un ExternalId ou d'un ExternalLink, puis autoriser l'insertion en base.

Cette chaîne est importante, car une importation directe depuis une source externe exposerait la plateforme à des données incomplètes ou incompatibles avec ses parcours internes. Le backend joue donc un rôle de filtre et de garant de cohérence avant toute insertion MySQL.

\subsection{Reverse flow de publication vers LinkedIn}

En complément de l'import d'offres externes, le projet intègre également un flux inverse de publication vers LinkedIn. Le principe retenu consiste à préparer dans le backend \textit{HireCue} les données nécessaires à la publication d'une offre, puis à utiliser l'API UniPile pour gérer le posting LinkedIn. Selon le niveau d'orchestration souhaité, \textit{n8n} peut intervenir comme couche de coordination entre le backend et UniPile afin d'enchaîner la préparation, la transmission et le suivi du flux sortant.

Ce reverse flow ne se limite pas à pousser un contenu ou une offre vers LinkedIn. Il prévoit également une tâche d'\textit{interaction tracking} après publication. Cette tâche consiste à suivre les commentaires, les réactions, les mentions J'aime et les statistiques d'engagement associés au post publié. Les informations récupérées peuvent être stockées dans le backend et dans la base de données afin de conserver l'identifiant du post, l'état de publication, les interactions observées et les indicateurs de performance. Cette logique transforme le flux sortant en boucle de suivi, car la plateforme ne se contente plus de publier : elle vérifie aussi l'impact du contenu et prépare son exploitation analytique.

Ce choix est plus adapté qu'un scraping navigateur pour un flux sortant, car il repose sur une couche API plus propre et plus défendable dans une architecture applicative. Il permet aussi de rapprocher la publication LinkedIn du tableau de bord recruteur ou administrateur, dans lequel les statistiques d'engagement pourront être consultées pour évaluer la visibilité d'une offre ou d'un contenu. Cette fonctionnalité doit toutefois être présentée avec mesure. Les éléments observés montrent surtout une base d'intégration et une logique d'orchestration préparée. Si toutes les étapes de publication et de suivi ne sont pas complètement stabilisées, le rapport doit les décrire comme des fonctionnalités intégrées ou en cours de consolidation, plutôt que comme un flux totalement industrialisé.

\subsection{Workflow d'automatisation des messages LinkedIn}

Le workflow d'automatisation des messages LinkedIn prolonge les scénarios de prospection de \textit{HireCue} vers des profils décisionnaires, notamment les CEO, founders ou owners. Son objectif est de structurer une prise de contact progressive, suivie et traçable, sans demander au recruteur de répéter manuellement chaque étape.

Le processus commence par la récupération des leads cibles. Ces leads peuvent provenir d'une recherche LinkedIn ou d'une liste déjà consolidée dans le backend. Avant l'envoi, le workflow vérifie leur statut afin d'éviter les doublons, les contacts déjà traités ou les profils qui ne correspondent pas aux critères définis. Lorsque le lead est éligible, \textit{n8n} orchestre l'envoi d'une demande de connexion LinkedIn via UniPile, puis met à jour le statut associé dans le backend.

La suite du workflow dépend de l'acceptation de la connexion. Le système peut attendre un retour ou détecter l'acceptation à travers un événement de synchronisation. Lorsque la connexion est acceptée, un message de suivi personnalisé est préparé puis envoyé automatiquement. Chaque étape donne lieu à une mise à jour de statut : lead identifié, invitation envoyée, connexion acceptée, message de suivi transmis ou action en attente. Cette conservation des traces permet d'assurer un suivi commercial plus rigoureux et d'alimenter ultérieurement des indicateurs analytiques tels que le nombre d'invitations envoyées, le taux d'acceptation ou le taux de réponse.

\section{Module décisionnel Power BI}

Cette section présente la réalisation de la couche décisionnelle intégrée au projet. Le module Power BI \cite{powerbi_docs} constitue la couche de Business Intelligence de \textit{HireCue}. Il complète les modules IA, les workflows \textit{n8n} et les automatisations par une lecture visuelle des données issues de MySQL. Son objectif n'est pas de remplacer les parcours applicatifs existants, mais de transformer les données opérationnelles de la plateforme en indicateurs lisibles pour les décideurs, les administrateurs et les responsables opérationnels.

La réalisation de cette partie suit une logique proche de CRISP-DM. La compréhension métier correspond à l'identification du besoin décisionnel : suivre l'activité globale de \textit{HireCue}, mesurer les volumes de recrutement, observer l'usage de l'IA et surveiller certains indicateurs opérationnels. La compréhension des données porte ensuite sur les tables MySQL disponibles, tandis que la préparation est assurée avec Power Query. Les mesures DAX et les visualisations permettent enfin de construire les KPIs, puis d'évaluer la lisibilité et l'utilité des dashboards produits.

L'intégration web des rapports dans l'application React / Node.js via une balise \texttt{iframe} relève davantage de la démarche Scrum. Elle concerne en effet une évolution applicative progressive : ajout d'une interface, intégration des vues Power BI, navigation entre dashboards et validation de l'affichage dans le navigateur.

\begin{table}[H]
\centering
\renewcommand{\arraystretch}{1.25}
\begin{tabular}{|p{3.2cm}|p{10.4cm}|}
\hline
\textbf{Technologie} & \textbf{Rôle dans le module} \\
\hline
MySQL & Source principale des données analytiques de \textit{HireCue}. \\
\hline
Power BI & Création des dashboards décisionnels et des visualisations. \\
\hline
Power Query & Nettoyage, transformation et préparation des données. \\
\hline
DAX & Création des mesures, calculs et KPIs. \\
\hline
React.js & Interface web permettant d'afficher et de naviguer entre les dashboards. \\
\hline
Node.js & Support applicatif/backend de l'intégration web. \\
\hline
\texttt{iframe} & Intégration des rapports Power BI dans l'interface web. \\
\hline
LSTM & Modèle Deep Learning utilisé pour la prévision de séries temporelles. \\
\hline
GRU & Modèle Deep Learning comparé à LSTM pour l'expérimentation prédictive. \\
\hline
\end{tabular}
\caption{Technologies utilisées dans le module décisionnel Power BI}
\label{tab:powerbi_technologies}
\end{table}

\subsection{Connexion de Power BI à la base MySQL de HireCue}

Power BI Desktop est connecté à la base MySQL \cite{mysql_docs} de \textit{HireCue}. Les données exploitées proviennent donc du système réel de la plateforme et non d'un fichier isolé ou d'un export statique. Cette connexion permet de construire une couche analytique cohérente avec les parcours candidat, recruteur et administrateur, tout en évitant la saisie manuelle de données dans les dashboards.

Selon les données disponibles dans la base, les tables exploitées peuvent inclure notamment des tables métier comme \texttt{users}, \texttt{joblist}, \texttt{candidate\_job} et \texttt{Events}. Elles peuvent aussi comprendre, lorsque leur présence est confirmée dans le modèle, des tables de suivi ou d'analyse telles que \texttt{userLogs}, \texttt{LLMConsumptions}, \texttt{subscriptions}, \texttt{score\_reports} ou des tables équivalentes liées aux scores, aux entretiens, aux workflows et aux interactions. Cette formulation reste volontairement prudente, car le modèle analytique doit suivre les tables réellement présentes et les besoins de reporting validés.

La logique générale peut être résumée par le flux suivant : MySQL \(\rightarrow\) Power BI Desktop \(\rightarrow\) Power Query \(\rightarrow\) modèle de données \(\rightarrow\) mesures DAX \(\rightarrow\) dashboards. Cette chaîne permet aux tableaux de bord de reposer directement sur les données opérationnelles de la plateforme, tout en conservant une étape de préparation avant la visualisation.

\begin{figure}[H]
    \centering
    \includegraphics[width=0.95\textwidth]{Coonectpowerbi-mysql.png}
    \caption{Connexion de Power BI à la base de données MySQL de HireCue}
    \label{fig:powerbi_mysql_connection}
\end{figure}

La figure \ref{fig:powerbi_mysql_connection} illustre le principe de connexion entre Power BI et la base MySQL. Les données récupérées alimentent ensuite les processus de nettoyage, de transformation, de modélisation et de création des indicateurs décisionnels. Cette architecture limite les manipulations manuelles et rapproche les dashboards de l'état réel de l'activité \textit{HireCue}.

\subsection{Nettoyage et transformation avec Power Query}

Après la connexion à MySQL, Power Query \cite{powerquery_docs} intervient comme phase de préparation des données avant la construction des visualisations. Cette étape correspond à la phase de préparation de CRISP-DM : elle transforme des données opérationnelles en données plus propres, plus lisibles et plus adaptées à l'analyse décisionnelle.

Les transformations réalisées peuvent comprendre le traitement des valeurs nulles, la correction des types de données, le formatage des dates, la normalisation des colonnes et le renommage de certains champs afin d'améliorer leur lisibilité dans les rapports. Les données non pertinentes pour l'analyse peuvent être filtrées, tandis que les colonnes nécessaires aux KPIs sont préparées ou vérifiées.

Dans le modèle préparé, des transformations simples mais importantes ont été utilisées. Les champs de date ont été convertis dans un format exploitable pour les axes temporels, les colonnes techniques ont été renommées pour devenir lisibles dans les visuels, les valeurs nulles ont été remplacées ou filtrées lorsqu'elles empêchaient le calcul d'un KPI, et certains champs numériques ont été typés correctement avant la création des mesures. Par exemple, les dates de création de jobs, les événements, les statuts de candidature et les consommations IA devaient être homogénéisés avant d'être utilisés dans les graphiques.

Power Query permet également de créer ou de vérifier les relations entre les tables utilisées dans le modèle. Ce point est essentiel lorsque les indicateurs combinent des informations issues des candidats, des offres, des événements, des logs, des abonnements ou des consommations IA. Cette préparation limite les incohérences dans les dashboards et garantit que les indicateurs ne reposent pas uniquement sur des données brutes, mais sur un modèle analytique structuré.

\subsection{Création des KPIs et mesures DAX}

La construction des dashboards repose ensuite sur la création de KPIs à partir des données MySQL préparées avec Power Query. DAX \cite{dax_docs} est utilisé pour définir les mesures analytiques nécessaires aux cartes KPI, aux graphiques, aux entonnoirs et aux comparaisons dynamiques. Ces mesures permettent aux indicateurs de réagir aux filtres, aux périodes et aux dimensions sélectionnées dans Power BI.

Les KPIs retenus ne reposent pas sur des chiffres ajoutés dans le rapport, mais sur les noms et familles d'indicateurs visibles dans les dashboards. Ils sont regroupés selon leur usage principal : pilotage stratégique, suivi IA / LLM, analyse du recrutement et supervision opérationnelle.

Parmi les mesures utilisées, \textit{Total Applications} correspond au volume de candidatures enregistrées, \textit{Active Recruiters} compte les recruteurs actifs dans la base, \textit{Jobs Created} suit le nombre d'offres créées et \textit{Completed Events} isole les événements terminés. Ces mesures restent volontairement simples dans leur principe, mais elles deviennent utiles lorsqu'elles sont combinées avec les filtres Power BI, les périodes, les entreprises, les pays ou les types d'emploi. D'autres mesures, comme les tokens IA ou les coûts par fournisseur, exploitent les tables de consommation afin de rendre l'usage des services IA plus observable.

\begin{table}[H]
\centering
\renewcommand{\arraystretch}{1.25}
\begin{tabular}{|p{3.2cm}|p{10.4cm}|}
\hline
\textbf{Groupe de KPIs} & \textbf{Indicateurs suivis} \\
\hline
KPIs stratégiques & \textit{Active Companies}, \textit{Active Recruiters}, \textit{Jobs Created}, \textit{Total Applications}, \textit{Active Subscriptions}, \textit{Revenue Proxy}, \textit{Selection Rate}. \\
\hline
KPIs IA / LLM & \textit{Total AI Tokens}, \textit{Average Score With AI}, \textit{LLM Calls}, \textit{Total AI Cost}, \textit{Avg Tokens Per Call}, \textit{Average AI Duration}, \textit{AI Cost by Provider}, \textit{AI Tokens Consumption by Task}, \textit{AI Tasks Distribution}. \\
\hline
KPIs recrutement & \textit{Recruitment Funnel}, \textit{Total Applications}, \textit{Video Interview Invites}, \textit{Selected Candidates}, \textit{Average Final Score}, \textit{Average Score Comparison}, \textit{Jobs by Status}, \textit{Jobs Growth Over Time}. \\
\hline
KPIs opérationnels & \textit{Events Overview}, \textit{Completed Events}, \textit{Planned Events}, \textit{Open Events}, \textit{Upcoming Today}, \textit{Recent Activity Logs}, \textit{Candidate Pipeline by Status}, \textit{Roadmaps Generated}, \textit{Not Completed Interviews}, \textit{Jobs by Job Type}. \\
\hline
\end{tabular}
\caption{Familles de KPIs exploitées dans les dashboards Power BI}
\label{tab:powerbi_kpis}
\end{table}

Les KPIs stratégiques aident les décideurs à suivre la performance générale de la plateforme. Les KPIs IA / LLM permettent de mesurer l'usage des services intelligents, notamment la consommation de tokens, les appels LLM, les coûts associés et les durées moyennes observées. Les KPIs de recrutement donnent une lecture du tunnel candidat, des candidatures, des sélections, des scores et de l'évolution des offres. Enfin, les KPIs opérationnels facilitent le suivi quotidien des événements, des logs, des statuts candidats, des roadmaps et des entretiens non complétés.

\subsection{Dashboard Executive Overview}

Le dashboard \textit{Executive Overview} est conçu comme une vue stratégique destinée aux décideurs et aux responsables qui ont besoin d'une lecture synthétique de l'activité globale de \textit{HireCue}. Il regroupe des cartes KPI, des graphiques temporels, des distributions et des indicateurs de performance afin de faciliter une lecture rapide des tendances principales.

Les cartes KPI visibles dans cette vue incluent notamment \textit{Active Companies}, \textit{Active Recruiters}, \textit{Jobs Created}, \textit{Total Applications}, \textit{Active Subscriptions}, \textit{Total AI Tokens}, \textit{Revenue Proxy} et \textit{Selection Rate}. Ces indicateurs donnent une première lecture de la taille de l'activité, de l'utilisation de la plateforme, de la dynamique des candidatures et de l'usage global de l'IA.

Le dashboard contient également des visualisations telles que \textit{Jobs Growth Over Time}, \textit{Revenue Proxy Over Time}, \textit{Subscription Distribution}, \textit{AI Cost Over Time}, \textit{Recruitment Funnel}, \textit{AI Tokens Consumption by Task}, \textit{Jobs by Status}, \textit{Average Score Comparison} et \textit{Completed Events}. L'ensemble permet de suivre la performance générale, d'observer les tendances de recrutement, de visualiser l'usage de l'IA et d'appuyer la prise de décision à partir d'indicateurs consolidés.

\begin{figure}[H]
    \centering
    \includegraphics[width=\textwidth]{Dashboard1.png}
    \caption{Dashboard Executive Overview de HireCue}
    \label{fig:dashboard-executive-overview}
\end{figure}

La figure \ref{fig:dashboard-executive-overview} présente la vue \textit{Executive Overview} intégrée dans le module décisionnel. Elle regroupe les indicateurs stratégiques liés aux entreprises actives, recruteurs actifs, offres créées, candidatures, abonnements, consommation de tokens IA, revenu proxy et taux de sélection. Cette vue offre une lecture synthétique de la performance globale de la plateforme.

\subsection{Dashboard Product \& Operations Analytics}

Le dashboard \textit{Product \& Operations Analytics} adopte une lecture plus opérationnelle. Il est destiné au suivi produit, aux workflows et à l'activité quotidienne de la plateforme. Là où \textit{Executive Overview} donne une vue synthétique pour la décision, ce second dashboard fournit une analyse plus détaillée des usages, des statuts et des comportements observés.

Les indicateurs visibles couvrent notamment \textit{Average Score With AI}, \textit{LLM Calls}, \textit{Total AI Cost}, \textit{Average Final Score}, \textit{Avg Tokens Per Call}, \textit{Average AI Duration}, \textit{Selection Rate} et \textit{Selected Candidates}. Ces mesures permettent de suivre à la fois la performance du recrutement et l'usage des modules IA.

La vue opérationnelle comprend aussi \textit{Jobs by Job Type}, \textit{AI Cost by Provider}, \textit{Recent Activity Logs}, \textit{AI Tokens Consumption Over Time}, \textit{AI Tasks Distribution}, \textit{Candidate Pipeline by Status}, \textit{Events Overview}, \textit{Roadmaps Generated} et \textit{Not Completed Interviews}. Elle aide ainsi à surveiller les logs récents, comprendre la consommation de tokens et le coût IA, suivre les événements, analyser les statuts candidats et compléter la lecture stratégique du dashboard \textit{Executive Overview}.

\begin{figure}[H]
    \centering
    \includegraphics[width=\textwidth]{Dashboard2.png}
    \caption{Dashboard Product \& Operations Analytics de HireCue}
    \label{fig:dashboard-product-operations}
\end{figure}

La figure \ref{fig:dashboard-product-operations} illustre la vue \textit{Product \& Operations Analytics}. Cette page permet de suivre les indicateurs opérationnels de la plateforme, notamment les appels LLM, le coût IA, le score moyen avec IA, les tokens consommés, les logs récents, les événements, les roadmaps générées, les statuts candidats et les entretiens non complétés.

\subsection{Intégration des dashboards dans une application React / Node.js}

Une interface web a été créée afin d'afficher les dashboards Power BI directement depuis le navigateur. Les rapports sont intégrés dans l'application React / Node.js à l'aide d'une balise \texttt{iframe}. React gère l'interface utilisateur, l'organisation des vues et la navigation entre les dashboards, tandis que Node.js peut servir de support applicatif ou backend selon la structure retenue pour l'application.



Cette solution permet d'afficher les rapports dans une interface unifiée et de naviguer entre \textit{Executive Overview} et \textit{Product \& Operations Analytics}, notamment à travers une sidebar dédiée. Elle simplifie le déploiement des tableaux de bord tout en conservant les capacités natives de filtrage, d'interaction et de visualisation offertes par Power BI.

\subsection{Expérimentation AI Workforce Forecasting}

L'\textit{AI Workforce Forecasting} constitue une expérimentation Deep Learning intégrée à l'interface web. Elle explore la possibilité de prévoir certaines tendances RH ou opérationnelles à partir des données disponibles dans la plateforme. Son objectif est d'évaluer l'intérêt d'une dimension prédictive en complément des dashboards Power BI, qui restent principalement orientés visualisation et suivi décisionnel.

\begin{figure}[H]
    \centering
    \includegraphics[width=\textwidth]{AIForeCasting1.png}
    \caption{Vue générale du module AI Workforce Forecasting}
    \label{fig:ai-workforce-forecasting-overview}
\end{figure}

La figure \ref{fig:ai-workforce-forecasting-overview} montre l'interface générale du module. Elle permet de lancer l'entraînement du modèle, d'afficher les prédictions principales et de présenter les KPIs prédits pour les 30 prochains jours. Les indicateurs suivis incluent notamment \textit{Predicted Jobs}, \textit{Predicted Applications}, \textit{Predicted AI Tokens}, \textit{Predicted AI Cost}, \textit{Predicted Revenue Proxy} et \textit{Best Model Used}.

\begin{figure}[H]
    \centering
    \includegraphics[width=\textwidth]{AIForeCasting2.png}
    \caption{Résultats de prédiction du module AI Workforce Forecasting}
    \label{fig:ai-workforce-forecasting-results}
\end{figure}

La figure \ref{fig:ai-workforce-forecasting-results} présente les résultats de prédiction associés à l'expérimentation. Le module compare deux modèles de séries temporelles, LSTM \cite{hochreiter_lstm} et GRU \cite{cho_gru}, afin d'observer leur comportement sur plusieurs indicateurs comme les jobs créés, les candidatures, le coût IA et le revenu proxy. Les courbes permettent de comparer les valeurs réelles avec les valeurs prédites.

La comparaison des modèles repose sur des métriques d'évaluation classiques : MAE, RMSE et MAPE. La MAE mesure l'erreur moyenne absolue, la RMSE pénalise davantage les grandes erreurs et la MAPE exprime l'erreur relative en pourcentage. Le meilleur modèle est sélectionné selon les résultats observés. Cette brique ne remplace toutefois pas les KPIs décisionnels existants et ne doit pas être présentée comme un module totalement industrialisé. À ce stade, il s'agit d'une preuve de concept analytique qui nécessite davantage de données, de validation et d'évaluation avant un usage décisionnel complet.

\section{Automatisation des workflows}

Cette section isole la partie automatisation afin de ne pas confondre les modules applicatifs avec les flux externes. Elle se concentre donc sur l'orchestration, le nettoyage, la déduplication, le suivi des statuts et la transmission des données entre \textit{HireCue}, \textit{n8n} et les services tiers.

\subsection{Workflows pour le sourcing candidat}

Le sourcing candidat repose sur une orchestration progressive. Le recruteur saisit un critère de recherche, puis le backend transmet les paramètres utiles à \textit{n8n}. Le workflow lance alors PhantomBuster, conserve les informations de suivi du run et prépare la récupération des profils.

Une phase de traitement séparée vérifie ensuite la fin de l'extraction, récupère les résultats, normalise les champs, supprime les doublons et renvoie au backend uniquement les profils exploitables. Cette séparation rend le flux plus lisible et évite qu'une collecte longue ou instable ne bloque directement l'expérience recruteur.

\subsection{Scraping LinkedIn des profils}

Le scraping LinkedIn fournit des données utiles au sourcing, mais ces données ne sont pas directement considérées comme fiables. Les profils, titres, entreprises, localisations ou emails professionnels disponibles doivent être interprétés comme des informations externes à vérifier.

Selon le scénario, les résultats peuvent transiter par un support intermédiaire ou être renvoyés vers le backend après traitement. L'enjeu principal reste la qualité du résultat final : seuls les profils suffisamment structurés et non dupliqués peuvent être affichés ou préparés pour une exploitation dans \textit{HireCue}.

\subsection{Pipeline d'import des offres d'emploi LinkedIn}

Le pipeline d'import des offres suit la même logique de contrôle. Les offres récupérées depuis des sources externes sont reformattées, puis mappées vers la structure interne de \textit{HireCue}. Le backend vérifie ensuite la présence des champs indispensables et l'absence de doublons, notamment à partir des identifiants ou des liens externes.

Ce mécanisme est important dans le projet, car il permet d'intégrer des offres dans \textit{HireCue} sans dupliquer des enregistrements déjà présents ni casser la cohérence du modèle métier. Dans une logique d'exploitation admin, ce workflow peut également être associé à des paramètres de recherche ou de sélection de source afin de piloter plus finement l'import des offres externes selon la plateforme ciblée.

\subsection{Workflow d'envoi d'emails}

Les workflows d'envoi d'emails concernent principalement les scénarios d'outreach. Ils peuvent lire des leads, filtrer les emails valides, vérifier les doublons, préparer un contenu personnalisé et déclencher l'envoi via Gmail. La valeur du workflow dépend surtout de la traçabilité : chaque lead doit conserver un statut clair, par exemple préparé, envoyé, répondu, ignoré ou en échec.

L'évolution autour du CEO Email Outreach Tracking renforce ce suivi. Les événements importants sont transmis au backend, persistés dans MySQL et affichés dans un dashboard administrateur. Un second workflow peut détecter les réponses via Gmail Trigger, puis mettre à jour l'historique. Cette logique reste volontairement cadrée : l'objectif n'est pas l'envoi massif, mais un pilotage par batch avec des statuts vérifiables.

\section{Difficultés rencontrées}

Au cours du développement des différentes fonctionnalités intégrées à la plateforme \textit{HireCue}, plusieurs difficultés techniques et organisationnelles ont été rencontrées. Ces contraintes ont nécessité la mise en place de solutions adaptées afin de garantir la stabilité, la performance et la cohérence globale du système.

\subsection{Intégration dans une plateforme existante}

L'une des principales difficultés concernait l'intégration de nouvelles fonctionnalités dans une plateforme déjà opérationnelle. Les modules développés devaient s'adapter à l'architecture frontend et backend existante sans perturber les fonctionnalités déjà en production.

Pour répondre à cette contrainte, une approche progressive a été adoptée consistant à analyser les composants existants, isoler les nouvelles responsabilités dans des services dédiés et intégrer les fonctionnalités de manière incrémentale.

\subsection{Validation des sorties des modèles d'IA}

Certaines fonctionnalités reposent sur des modèles d'intelligence artificielle produisant des réponses au format JSON. Une réponse correcte sur le plan textuel peut néanmoins être inutilisable si sa structure ne respecte pas le format attendu par le backend.

Afin d'assurer la fiabilité des traitements, plusieurs mécanismes ont été mis en place : prompts contraints, validation de schéma, régénération automatique en cas d'erreur et fallback déterministe lorsque la réponse ne respecte pas les exigences minimales.

\subsection{Gestion des coûts et des quotas IA}

L'utilisation de services d'intelligence artificielle implique des contraintes liées au temps de calcul et à la consommation de ressources.

Pour limiter les appels inutiles aux modèles, un système de cache, de snapshots et de journalisation des consommations a été mis en œuvre. Cette approche permet de réutiliser les résultats déjà générés et d'améliorer les performances globales de la plateforme.

\subsection{Fiabilité des workflows automatisés}

Les workflows \textit{n8n} et certaines intégrations externes ont présenté plusieurs défis, notamment la variabilité des données récupérées, les erreurs d'exécution, les limitations imposées par certaines plateformes et l'hétérogénéité des formats de données.

Pour améliorer la robustesse de ces traitements, les workflows ont été isolés dans des processus dédiés intégrant des mécanismes de normalisation, de déduplication et de contrôle d'état.

\subsection{Choix des outils d'automatisation LinkedIn}

Le projet a nécessité l'utilisation de plusieurs outils spécialisés pour automatiser les interactions avec LinkedIn.

\textbf{PhantomBuster} s'est révélé particulièrement adapté aux tâches de scraping, de sourcing de candidats, d'extraction de profils et d'enrichissement de données. Il a notamment été utilisé pour la récupération de profils LinkedIn, le scraping de CEO et l'import d'offres externes.

À l'inverse, \textbf{UniPile} est davantage orienté communication et intégration API. Il facilite la gestion des messages LinkedIn, le suivi des conversations ainsi que la publication de contenus.

Cette complémentarité a conduit à adopter une architecture hybride où chaque outil est utilisé selon ses points forts plutôt que comme une solution universelle.

\subsection{Performance et traitements asynchrones}

Certaines fonctionnalités, telles que les recommandations intelligentes ou les campagnes de réengagement, nécessitent des traitements relativement coûteux.

Une exécution entièrement synchrone aurait dégradé l'expérience utilisateur. Pour éviter ce problème, une architecture asynchrone reposant sur RabbitMQ, des tâches en arrière-plan et un mécanisme de polling côté frontend a été mise en place.

Cette approche permet de conserver une interface réactive tout en exécutant des calculs plus complexes en arrière-plan.

\subsection{Cohérence entre le frontend et le backend}

Enfin, le maintien de la cohérence entre les services backend et les interfaces utilisateur a constitué un défi permanent.

Chaque fonctionnalité devait garantir une gestion homogène des statuts, des formats de réponse et des erreurs afin d'assurer une expérience utilisateur fluide et compréhensible. Ce travail d'alignement a accompagné l'ensemble des développements réalisés durant le stage.

\section{Conclusion}

Ce chapitre a présenté l'architecture générale de \textit{HireCue}, la lecture UML du système, la conception fonctionnelle et technique des modules, puis l'implémentation des fonctionnalités et des workflows les plus significatifs du projet. Il a aussi montré que les briques IA et les automatisations externes n'ont de valeur que lorsqu'elles restent reliées à des données métier fiables, à des contrôles techniques explicites et à des parcours utilisateur compréhensibles.

Le chapitre suivant pourra ainsi se concentrer sur la validation, les tests et les résultats obtenus à partir de ces modules et de leur intégration dans la plateforme.
